\documentclass[aspectratio=169]{beamer}

\usetheme{Madrid}
\usecolortheme{default}

\usepackage{amsmath,amssymb,amsfonts,bm}
\usepackage{mathtools}

\title{Linear Autoencoders Recover PCA}
\author{Morten Hjorth-Jensen}
\date{Spring 2026}

\begin{document}

\frame{\titlepage}

\begin{frame}{Statement of the result}
Let \(X\in \mathbb{R}^{d\times N}\) be a centered data matrix, whose columns are data vectors
\[
x_1,\dots,x_N\in \mathbb{R}^d,
\qquad
\frac1N \sum_{n=1}^N x_n = 0.
\]

Consider a linear autoencoder with encoder matrix \(W_e\in \mathbb{R}^{k\times d}\)
and decoder matrix \(W_d\in \mathbb{R}^{d\times k}\), where \(k<d\).

The reconstruction is
\[
\widehat X = W_d W_e X.
\]

We minimize the squared reconstruction error
\[
\mathcal L(W_d,W_e)
=
\|X-W_dW_eX\|_F^2.
\]

\medskip

\textbf{Claim:} every global minimizer spans the same \(k\)-dimensional principal subspace as PCA.
\end{frame}

\begin{frame}{What exactly must be proved?}
There are two levels of the result:

\begin{enumerate}
\item The optimal product
\[
A^\star = W_d^\star W_e^\star
\]
is the orthogonal projector onto the principal \(k\)-dimensional subspace.

\item The encoder and decoder individually need not equal the PCA matrices uniquely, but their column/row spaces recover the same subspace.
\end{enumerate}

So the precise statement is not that
\[
W_e^\star = U_k^T,\qquad W_d^\star = U_k
\]
must hold uniquely, but that the optimal autoencoder reconstructs by projection onto the PCA subspace.
\end{frame}

\begin{frame}{Step 1: Reformulate the problem}
Define
\[
A = W_d W_e \in \mathbb{R}^{d\times d}.
\]

Because \(W_d\in \mathbb{R}^{d\times k}\) and \(W_e\in \mathbb{R}^{k\times d}\),
the matrix \(A\) satisfies
\[
\operatorname{rank}(A)\le k.
\]

Thus the linear autoencoder optimization can be rewritten as
\[
\min_{W_d,W_e}\|X-W_dW_eX\|_F^2
\quad \Longrightarrow \quad
\min_{\operatorname{rank}(A)\le k}\|X-AX\|_F^2.
\]

So the problem is to find the best rank-\(k\) linear operator \(A\) that approximates the identity on the dataset.
\end{frame}

\begin{frame}{Step 2: Introduce the sample covariance matrix}
Since the data are centered, the empirical covariance matrix is
\[
S = \frac1N XX^T.
\]

It is symmetric and positive semidefinite, so it admits the eigendecomposition
\[
S = U \Lambda U^T,
\]
where
\[
U = (u_1,\dots,u_d),\qquad
\Lambda=\operatorname{diag}(\lambda_1,\dots,\lambda_d),
\]
with
\[
\lambda_1\ge \lambda_2\ge \dots \ge \lambda_d \ge 0.
\]

The PCA subspace of dimension \(k\) is
\[
\mathcal U_k = \operatorname{span}\{u_1,\dots,u_k\}.
\]
\end{frame}

\begin{frame}{Step 3: Expand the loss}
We compute
\[
\|X-AX\|_F^2
=
\operatorname{Tr}\Big((X-AX)^T(X-AX)\Big).
\]

Expanding,
\[
\|X-AX\|_F^2
=
\operatorname{Tr}(X^TX)
-2\operatorname{Tr}(X^TAX)
+\operatorname{Tr}(X^TA^TAX).
\]

Using cyclicity of the trace,
\[
\operatorname{Tr}(X^TAX)=\operatorname{Tr}(AXX^T),
\qquad
\operatorname{Tr}(X^TA^TAX)=\operatorname{Tr}(A^TAXX^T).
\]

Since \(XX^T = N S\), we get
\[
\|X-AX\|_F^2
=
N\Big[
\operatorname{Tr}(S)
-2\operatorname{Tr}(AS)
+\operatorname{Tr}(A^TAS)
\Big].
\]
\end{frame}

\begin{frame}{Step 4: Restrict first to orthogonal projectors}
Suppose \(P\) is an orthogonal projector of rank \(k\):
\[
P^2=P,\qquad P^T=P,\qquad \operatorname{rank}(P)=k.
\]

Then
\[
P^TP = P,
\]
and the loss becomes
\[
\|X-PX\|_F^2
=
N\Big[\operatorname{Tr}(S)-2\operatorname{Tr}(PS)+\operatorname{Tr}(PS)\Big]
=
N\Big[\operatorname{Tr}(S)-\operatorname{Tr}(PS)\Big].
\]

Therefore minimizing the loss over rank-\(k\) orthogonal projectors is equivalent to maximizing
\[
\operatorname{Tr}(PS).
\]
\end{frame}

\begin{frame}{Step 5: Parameterize a rank-\(k\) projector}
Any rank-\(k\) orthogonal projector can be written as
\[
P = BB^T,
\]
where \(B\in \mathbb{R}^{d\times k}\) has orthonormal columns:
\[
B^TB=I_k.
\]

Then
\[
\operatorname{Tr}(PS)
=
\operatorname{Tr}(BB^TS)
=
\operatorname{Tr}(B^TSB).
\]

So the optimization becomes
\[
\max_{B^TB=I_k}\operatorname{Tr}(B^TSB).
\]
This is the standard variational characterization of PCA.
\end{frame}

\begin{frame}{Step 6: Use the eigendecomposition of \(S\)}
Write
\[
S = U\Lambda U^T.
\]

Let
\[
C = U^T B.
\]

Since \(U\) is orthogonal and \(B^TB=I_k\),
\[
C^TC = B^TUU^TB = B^TB = I_k.
\]

Now
\[
\operatorname{Tr}(B^TSB)
=
\operatorname{Tr}(B^TU\Lambda U^TB)
=
\operatorname{Tr}(C^T\Lambda C).
\]
\end{frame}

\begin{frame}{Step 7: Expand the trace in eigen-coordinates}
Let \(c_{ij}\) be the entries of \(C\). Then
\[
\operatorname{Tr}(C^T\Lambda C)
=
\sum_{j=1}^k \sum_{i=1}^d \lambda_i c_{ij}^2
=
\sum_{i=1}^d \lambda_i \sum_{j=1}^k c_{ij}^2.
\]

Define
\[
\alpha_i = \sum_{j=1}^k c_{ij}^2.
\]

Then
\[
\operatorname{Tr}(C^T\Lambda C)=\sum_{i=1}^d \lambda_i \alpha_i.
\]

Because \(C^TC=I_k\),
\[
\sum_{i=1}^d \alpha_i
=
\sum_{i=1}^d \sum_{j=1}^k c_{ij}^2
=
\sum_{j=1}^k \sum_{i=1}^d c_{ij}^2
=
k.
\]

Also each \(\alpha_i\) satisfies
\[
0\le \alpha_i \le 1.
\]
\end{frame}

\begin{frame}{Step 8: Maximize the weighted sum}
We must maximize
\[
\sum_{i=1}^d \lambda_i \alpha_i
\]
subject to
\[
0\le \alpha_i \le 1,
\qquad
\sum_{i=1}^d \alpha_i = k.
\]

Since the eigenvalues are ordered
\[
\lambda_1\ge \lambda_2\ge \cdots \ge \lambda_d,
\]
the maximum is obtained by putting all the weight on the largest \(k\) eigenvalues:
\[
\alpha_1=\cdots=\alpha_k=1,
\qquad
\alpha_{k+1}=\cdots=\alpha_d=0.
\]

Hence
\[
\max_{B^TB=I_k}\operatorname{Tr}(B^TSB)
=
\sum_{i=1}^k \lambda_i.
\]

This is achieved when the columns of \(B\) span the top-\(k\) eigenspace:
\[
B = U_k R,
\]
where \(U_k=(u_1,\dots,u_k)\) and \(R\in \mathbb{R}^{k\times k}\) is orthogonal.
\end{frame}

\begin{frame}{Conclusion for orthogonal projectors}
Thus the optimal orthogonal projector is
\[
P^\star = U_k U_k^T.
\]

The corresponding reconstruction error is
\[
\|X-P^\star X\|_F^2
=
N\left(\operatorname{Tr}(S)-\sum_{i=1}^k \lambda_i\right)
=
N\sum_{i=k+1}^d \lambda_i.
\]

This is exactly the PCA reconstruction error.
\end{frame}

\begin{frame}{Step 9: Why this also solves the full linear autoencoder problem}
So far we optimized over rank-\(k\) orthogonal projectors.
But the autoencoder problem allows any factorization
\[
A = W_dW_e
\]
with \(\operatorname{rank}(A)\le k\), not necessarily a projector.

We now show that no such \(A\) can do better than the PCA projector.
\end{frame}

\begin{frame}{Image-space argument}
Let
\[
\mathcal M = \operatorname{Im}(A)
\]
be the image of \(A\). Since \(\operatorname{rank}(A)\le k\), this is a subspace of dimension at most \(k\).

For each data vector \(x_n\), the reconstruction \(Ax_n\) lies in \(\mathcal M\).
Therefore \(Ax_n\) is some point in \(\mathcal M\), but not necessarily the \emph{orthogonal projection} of \(x_n\) onto \(\mathcal M\).

Let \(P_{\mathcal M}\) denote the orthogonal projector onto \(\mathcal M\).
By the defining property of orthogonal projection,
\[
\|x_n - P_{\mathcal M}x_n\|_2^2 \le \|x_n - Ax_n\|_2^2
\]
for every \(n\).
\]
Summing over \(n\),
\[
\|X-P_{\mathcal M}X\|_F^2 \le \|X-AX\|_F^2.
\]
\end{frame}

\begin{frame}{Finish the argument}
This means: for every rank-\(k\) matrix \(A\), there exists an orthogonal projector \(P_{\mathcal M}\) of rank at most \(k\) with no larger error.

Hence the global minimizer of
\[
\min_{\operatorname{rank}(A)\le k}\|X-AX\|_F^2
\]
may be taken to be an orthogonal projector.

But we already proved that the best rank-\(k\) orthogonal projector is
\[
P^\star = U_kU_k^T.
\]

Therefore
\[
\min_{\operatorname{rank}(A)\le k}\|X-AX\|_F^2
=
\|X-U_kU_k^TX\|_F^2.
\]

So the optimal linear autoencoder reconstructs by projection onto the PCA subspace.
\end{frame}

\begin{frame}{How to choose encoder and decoder}
A convenient choice is
\[
W_d^\star = U_k,
\qquad
W_e^\star = U_k^T.
\]

Then
\[
W_d^\star W_e^\star = U_kU_k^T = P^\star.
\]

So the encoder computes principal coordinates
\[
z = U_k^T x,
\]
and the decoder reconstructs by
\[
\hat x = U_k z = U_kU_k^T x.
\]

This is exactly PCA reconstruction.
\end{frame}

\begin{frame}{Non-uniqueness of the factorization}
The optimal product \(A^\star = U_kU_k^T\) is unique if the PCA subspace is unique, but the factors are not.

For any invertible matrix \(M\in \mathbb{R}^{k\times k}\),
\[
W_d = U_k M,
\qquad
W_e = M^{-1} U_k^T
\]
also satisfy
\[
W_dW_e = U_kU_k^T.
\]

Thus the learned latent coordinates are not unique:
they are determined only up to an invertible linear change of basis.
\end{frame}

\begin{frame}{Equivalent statement using SVD}
If
\[
X = U \Sigma V^T
\]
is the singular value decomposition of \(X\), then
\[
S = \frac1N XX^T = \frac1N U\Sigma^2 U^T.
\]

Thus the principal directions are the left singular vectors of \(X\).

The optimal linear autoencoder reconstructs with
\[
A^\star = U_kU_k^T,
\]
where \(U_k\) contains the first \(k\) left singular vectors.

So the result is equivalent to the rank-\(k\) approximation theorem of Eckart--Young--Mirsky.
\end{frame}

\begin{frame}{Final theorem}
\textbf{Theorem.}
Let \(X\in \mathbb{R}^{d\times N}\) be centered data and let \(k<d\).
Consider the linear autoencoder objective
\[
\min_{W_d\in\mathbb{R}^{d\times k},\; W_e\in\mathbb{R}^{k\times d}}
\|X-W_dW_eX\|_F^2.
\]

Then every global minimizer satisfies
\[
\operatorname{Im}(W_d^\star)=\mathcal U_k,
\]
the principal \(k\)-dimensional PCA subspace, and the optimal reconstruction map is
\[
W_d^\star W_e^\star = U_kU_k^T.
\]

Hence linear autoencoders recover PCA.
\end{frame}

\begin{frame}{Summary}
Main steps of the proof:
\begin{enumerate}
\item Rewrite the autoencoder as a rank-constrained linear map \(A=W_dW_e\).
\item Show that for fixed image subspace, orthogonal projection is optimal.
\item Reduce the optimization to maximizing
\[
\operatorname{Tr}(B^TSB)
\quad \text{subject to} \quad B^TB=I_k.
\]
\item Use the eigendecomposition of \(S\) to show the maximizer is the span of the top \(k\) eigenvectors.
\end{enumerate}

Therefore the optimal linear autoencoder performs exactly PCA projection and reconstruction.
\end{frame}

\end{document}
